\documentclass[a4paper,12pt,english]{article}
\usepackage{glossaries}
\usepackage[T1]{fontenc}
\usepackage{babel}
\usepackage{graphicx}
\usepackage[table,xcdraw]{xcolor}
\usepackage{hyperref}
\usepackage{blindtext}
\usepackage{geometry}
\usepackage{parskip}
\usepackage{mathtools}
\usepackage{siunitx}
\usepackage{listings}
\usepackage{csquotes}
\usepackage{caption}
\usepackage{subcaption}
\usepackage{comment}
\usepackage{pdfpages}
\usepackage{pict2e}

% \DeclareRobustCommand{\slashcirc}{{\mathpalette\doslashcirc\relax}}

% \makeatletter
% \newcommand\doslashcirc[2]{%
%   \sbox\z@{$#1\m@th\circ$}%
%   \setlength\unitlength{\wd\z@}
%   \begin{picture}(1,1)
%   \roundcap
%   \put(0,0){\box\z@}
%   \put(0,0){\line(1,1){1}}
%   \end{picture}%
% }
% \makeatother


% %% Some packages you will need
% \usepackage{pgfplots}
% \usepackage{pgfplotstable}
% \usepackage{booktabs}
% \usepackage{array}
% \usepackage{colortbl}


\definecolor{arduinoorange}{HTML}{FFA500}
\definecolor{arduinogray}{HTML}{808080}
\definecolor{arduinoblue}{HTML}{007ACC}
\definecolor{arduinogreen}{HTML}{469B00}

\lstset{
  language=C++,
  basicstyle=\ttfamily\footnotesize,
  keywordstyle=\color{arduinoorange},
  stringstyle=\color{arduinogreen},
  commentstyle=\color{arduinogray},
  moredelim=[s][\color{arduinoblue}]{\#}{\ },
  morekeywords={digitalRead,digitalWrite,pinMode,analogRead,analogWrite,Serial,begin,HIGH,LOW},
  frame=tb,
  tabsize=4,
  showstringspaces=false,
  breaklines=true,
  numbers=left,
  numberstyle=\tiny\color{arduinogray},
  numbersep=5pt,
  extendedchars=true,
  literate={á}{{\'a}}1 {ã}{{\~a}}1 {é}{{\'e}}1,
}

\lstdefinestyle{Arduino}
{
  language=C++,
  basicstyle=\ttfamily\footnotesize,
  keywordstyle=\color{arduinoorange},
  stringstyle=\color{arduinogreen},
  commentstyle=\color{arduinogray},
  moredelim=[s][\color{arduinoblue}]{\#}{\ },
  morekeywords={digitalRead,digitalWrite,pinMode,analogRead,analogWrite,Serial,begin},
  frame=tb,
  tabsize=4,
  showstringspaces=false,
  breaklines=true,
  numbers=left,
  numberstyle=\tiny\color{arduinogray},
  numbersep=5pt,
  extendedchars=true,
  literate={á}{{\'a}}1 {ã}{{\~a}}1 {é}{{\'e}}1,
  backgroundcolor=\color{black!85},
  rulecolor=\color{arduinoorange},
  frame=single,
  frameround=tttt,
  framexleftmargin=6pt,
  framexrightmargin=6pt,
  framextopmargin=6pt,
  framexbottommargin=6pt,
  breaklines=true,
  postbreak=\raisebox{0ex}[0ex][0ex]{\ensuremath{\color{red}\hookrightarrow\space}},
}

\usepackage[
    backend=biber,
    backref=true,
    backrefstyle=none,
    sortcites=true,
    sorting=none,
    doi=false, % doi informatie wordt niet weergegeven
    %uniquename=true,
    %uniquelist=true,
    maxcitenames=3,
    %issn=false, werkt niet
    language=american
]{biblatex}
\addbibresource{information/Sources.bib}
\DefineBibliographyStrings{english}{
    backrefpage = {blz.},
    backrefpages = {blz.},
}
\makeglossaries
\definecolor{Grey1}{HTML}{343434}
\graphicspath{{./Media/Figuren/}}
 \geometry{
 a4paper,
 total={170mm,257mm},
 left=20mm,
 top=20mm,
 }
\hypersetup{
    colorlinks=true,
    linkcolor=blue,
    filecolor=magenta,      
    urlcolor=cyan,
    pdftitle={Overleaf Example},
    pdfpagemode=FullScreen,
    }


\begin{document}
\title{
\includegraphics[width=3.5in]{img/Logo.png} \\
\vspace*{1in}
\textbf{Project 4}\\
\textit{Plan of Approach}\\
Version 1
}
\author{
\vspace*{0.5in} \\
  Written by:\\
  Laurens van der Drift\\
  Justin van der Reijden\\
  Luuk van Kappel\\
  Marnix Harmsen\\
		\vspace*{0.2in} \\
		The Hague University of Applied Sciences\\
        \textbf{Electrical Engineering}\\
        Delft, The Netherlands
       } 
\maketitle
\phantomsection
\section*{Version History} \addcontentsline{toc}{section}{Version History}

\begin{table}[h]
\begin{tabular}{|l|l|l|l|}
\hline
\rowcolor[HTML]{4472C4} 
{\color[HTML]{FFFFFF} \textbf{Version}} &
  {\color[HTML]{FFFFFF} \textbf{Date}} &
  {\color[HTML]{FFFFFF} \textbf{Changes}} &
  {\color[HTML]{FFFFFF} \textbf{Author}} \\ \hline
\rowcolor[HTML]{D9E1F2} 
1.0 &
  \multicolumn{1}{c|}{\cellcolor[HTML]{D9E1F2}05-02-2024} &
 N.v.t. &
  Alset Innovations \\ \hline

% \rowcolor[HTML]{FFFFFF} 
% 2.0 &
%   \multicolumn{1}{c|}{\cellcolor[HTML]{FFFFFF}7-4-2023} &
%  Feedback van FeedbackFruits toegepast &
%   Infra   Vroom \\ \hline

% \rowcolor[HTML]{D9E1F2} 
% 3.0 &
%   \multicolumn{1}{c|}{\cellcolor[HTML]{D9E1F2}4-6-2023} &
%  Bijgewerkt voor Assessment 3 &
%   Infra   Vroom \\ \hline

\end{tabular}
\end{table}
% %\addcontentsline{toc}{section}{Verklarende Woordenlijst}
\printglossaries
\newglossaryentry{tender}
{
    name=\textit{tender},
    description={De inschrijving om op een kavel een wind turbine park te bouwen.}
}

% \phantomsection
\addcontentsline{toc}{section}{List of Figures}
\listoffigures

\newpage
\tableofcontents
\section{The Assignment}
In the ambitious realm of the \textbf{PROJE4} course, \textbf{Alset Innovations} is tasked with the exciting challenge of designing and implementing a cutting-edge solution for Smart Energy Devices. The focal point of our project revolves around creating a sensored brushless-DC (BLDC) motor controller for an electric bicycle. This endeavor aligns seamlessly with the broader objective of utilizing intelligent embedded technology to enhance efficiency and conserve energy in various applications.

\subsection{Project Scope:}
Our primary goal is to develop a motor controller capable of seamlessly driving a \SI{24}{\volt}, \SI{78}{\watt} BLDC motor in a closed-loop configuration using hall sensor feedback. Our innovative twist lies in the selection of the \textbf{STM32F103C8T6} microcontroller, a new and unexplored avenue, which distinguishes us from our peers who may opt for the more conventional \textbf{RP2040}. This strategic decision places \textbf{Alset Innovations} at the forefront of innovation of this course.

% \subsection{Key Technical Requirements:}
% \begin{enumerate}
%     \item Utilize discrete transistors to control the three motor phases.
%     \item Design a custom printed circuit board (PCB) for the entire system, showcasing our prowess in hardware development.
%     \item Measure and display the current delivered to the motor through the user interface.
%     \item Implement a user-friendly control input (throttle) on the user interface for speed control.
%     \item Integrate the \textbf{STM32F103C8T6} microcontroller to process hall sensor signals, demonstrating our proficiency in microcontroller programming.
%     \item Deliver a comprehensive technical report detailing our design choices, simulations, schematics, and testing results.
% \end{enumerate}

\subsection{Strategic Approach:}
\textbf{Alset Innovations} recognizes the uniqueness of our project by opting for the \textbf{STM32} microcontroller, thereby charting a pioneering course in the implementation of smart energy devices. Our approach involves meticulous planning and execution, encompassing the following key phases:

\begin{enumerate}
    \item \textbf{Initial Research and Exploration:}
        \begin{itemize}
            \item Dive into the specifics of BLDC motor control, \textbf{STM32} microcontroller capabilities, and related electronic circuit theory.
            \item Explore the feasibility of our chosen microcontroller and its potential advantages in comparison to the \textbf{RP2040}.
        \end{itemize}

    \item \textbf{Hardware and PCB Design:}
        \begin{itemize}
            \item Research, compare and test different \textbf{MOSFETS} and \textbf{GATE} drivers.
            \item Develop a detailed schematic and layout for the PCB, focusing on optimal integration and functionality.
            \item Address the challenges posed by components like the \textbf{QFN-56} package of the \textbf{STM32} microcontroller.
        \end{itemize}

    \item \textbf{Software Development:}
        \begin{itemize}
            \item Craft robust code for the microcontroller, ensuring seamless motor control and user interface interaction.
            \item Conduct thorough testing on real prototype hardware, substantiating our software's reliability.
        \end{itemize}

    \item \textbf{Documentation and Reporting:}
        \begin{itemize}
            \item Maintain an exhaustive record of our design process, simulations, and results.
            \item Adhere to the specified template for the final technical report, showcasing our commitment to clarity and precision.
        \end{itemize}
\end{enumerate}

In summary, \textbf{Alset Innovations} stands poised to redefine the landscape of smart energy devices through our unique blend of innovation, strategic decision-making, and meticulous execution. Our Plan of Approach sets the stage for a transformative journey in the world of electrical engineering.

\section{Project Activity}
During this project a number of deadlines must be met, four in total. To meet these deadlines \textbf{Alset Innovations} has compiled a list of activities which must be met to complete the \textbf{PROJE4} course.

\subsection{Deadlines}
The first deadline for the \textbf{PROJE4} course is the completion of the Plan of Approach this must be completed and delivered on the 18th of February 23:59. After this \textbf{Alset Innovations} will be dedicated to completing assessment 1 and 2. For which the main focus will be on designing the code for the \textbf{STM32} as this is required for the first assessment. The secondary focus will be on designing a schematic for the motor driver which will be used to order parts and for assessment 2 to design a \textbf{PCB}. These assessments must be completed in week 7 and week 12 respectively.

The final assessment will be held in the 19th week of the project, for this assessment a presentation and a demonstration must be held. In addition to this a written report must also be delivered.

\subsection{Activities}
To meet these deadlines the following activities will be completed by \textbf{Alset Innovations}

\begin{enumerate}
    \item Complete planning to complete the Plan of Approach.
    \item Complete the Plan of Approach.
    \item Design a circuit to read the HAL sensors.
    \item Use the \textbf{STM32} to read the HAL sensors.
    \item Design code to control the timing of the \textbf{MOSFETS} using the HAL sensor data.
    \item Design code to create a PWM signal to control the motor RPM using an input as reference.
    \item Design code to create an interface
    \item Design a schematic with the most important components and order these parts.
    \item Test the function of the schematic mentioned above.
    \item Based on the testing add necessary components, protective circuitry and additional filters to the schematic. Order the parts that were added.
    \item Test the function of the updated schematic and continue to improve the schematic and test the schematic until the schematic functions as is required.
    \item After completing the schematic the PCB can be designed.
    \item 
\end{enumerate}
\section{Project Boundaries}
\section{Costs and revenues}
In pursuit of our Smart Energy Devices project, managing costs and assessing potential revenues are critical aspects. This section outlines the financial considerations associated with the development and implementation of our BLDC sensored motor controller.

\subsection{Budget Allocation}
Our project operates within a budget of 100 euros per group, excluding the cost of a reasonably sized bare PCB. This allocation covers various expenses, including components, tools, and potential unforeseen costs during the project's lifecycle.

\subsubsection{Component Orders}
All component orders are made through approved suppliers listed on Brightspace, with Conrad or Farnell being highly recommended due to their advantageous shipping policies. To streamline the ordering process, groups are required to follow these steps:

\begin{enumerate}
    \item Create a comprehensive component list in MS Excel using the provided template on Brightspace.
    \item Rename the file as \textbf{PROJE4 2023 Component Order List Group Number.xlsx}, where \textbf{Number} corresponds to the group's assigned number in Brightspace.
    \item Double-check the suitability of the component packages for hand soldering.
    \item Verify component availability and delivery times.
    \item Submit the order list to Zoja Donné (\textbf{Z.Donne@hhs.nl}), with the course organizer (\textbf{Stephen O’Loughlin S.D.Oloughlin@hhs.nl}) in CC.
    \item Regularly check with Zoja regarding delivery times until components are received.
\end{enumerate}

\subsection{Expense Tracking}
To ensure transparency and accountability, each group must diligently track all purchased items and include a detailed record in the final report.

\subsection{Revenue Projection}
While our primary focus is on academic and technical achievements, the potential for future revenue generation should not be overlooked. The skills developed during this project, coupled with the innovative nature of our sensored BLDC motor controller, may open opportunities for collaboration, licensing, or even commercialization.

\subsection{Cost-Benefit Analysis}
A comprehensive cost-benefit analysis will be conducted at the project's conclusion, considering both tangible and intangible factors. This analysis will provide insights into the efficiency of resource utilization, the project's overall viability, and potential areas for improvement in future iterations.

\subsection{Financial Planning}
Effective financial planning is essential for successful execution of the project. Regular budget reviews, coupled with prudent decision making regarding expenditures, will contribute to the overall financial health of the project.

In conclusion, the prudent management of costs and a forward-looking perspective on potential revenues are integral components of our Smart Energy Devices project. This financial approach ensures the sustainability of the project, fosters responsible resource utilization, and sets the stage for future innovations in the realm of smart energy devices.

\section{Organisation}
\section{Planning}





In this report, we outline the project plan and scheduling details for the development of Alset Innovations. The planning and execution of the project activities will be facilitated using GitHub's built-in roadmap feature. 

\subsection*{Project Timeline}

The project is segmented into three assessments, each with specific tasks and deadlines:

\begin{enumerate}
    \item \textbf{Plan of Approach:}February 5th to February 12th
    \item \textbf{Assessment 1:} February 13th to March 25th
    \item \textbf{Assessment 2:} March 26th to April 29th
    \item \textbf{Final Assessment:} April 30th to June 17th
\end{enumerate}

\subsection*{Project Activities and Deadlines}

\begin{enumerate}
    \item \textbf{Plan of Approach :}
    \begin{itemize}
        \item Complete the Plan of Approach document.
    \end{itemize}

    \item \textbf{Assessment 1 Tasks:}
    \begin{itemize}
        \item Design a circuit to read the HAL sensors.
        \item Develop code for reading the HAL sensors.
        \item Design code to control MOSFETS timing based on HAL sensor data.
        \item Develop code for interfacing with a touch-display or OLED display.
        \item Design a schematic for motor driving components.
        \item Document detailed software design requirements.
        \item Prepare well-documented code.
        \item Conduct tests and document results.
        \item Deliver a presentation.
    \end{itemize}

    \item \textbf{Assessment 2 Tasks:}
    \begin{itemize}
        \item Enhance schematic based on testing.
        \item Design and test PCB.
        \item Document detailed hardware design requirements.
        \item Provide PCB design files and datasheets.
        \item Create a clear PDF showcasing schematic and PCB layout.
        \item Summarize project progress and test results.
        \item Deliver a presentation.
    \end{itemize}

    \item \textbf{Final Assessment Tasks:}
    \begin{itemize}
        \item Deliver a final presentation.
        \item Conduct a project demonstration.
        \item Submit the final project report.
    \end{itemize}
\end{enumerate}

\subsection*{Resource Allocation}

Project members are expected to allocate a minimum of 8 hours per week towards project activities to ensure completion of tasks.

\subsection*{Project Management Tool}

GitHub will be utilized for project planning, task tracking, and collaboration among team members. The built-in roadmap feature will be employed to visualize project progress and milestones.





% Uit dit onderdeel moet de lezer begrijpen wat je allemaal gaat doen om het project succesvol af te
% ronden. Je geeft zo gedetailleerd mogelijk aan welke taken (activiteiten) er tijdens het project
% uitgevoerd moeten worden.
% - Achterhaal ALLE activiteiten die moeten worden uitgevoerd voor het project. Het is anders
% onmogelijk om een goede planning en kostenberekening te maken. Een goed overzicht van alle
% activiteiten is al een halve planning.
% - Ook de voorbereiding ook op in de activiteitenlijst denk hierbij aan:
% o Verzamelen en bestuderen informatie;
% o Gesprekken met de opdrachtgever;
% o Maken van een (concept) plan van aanpak;
% o ….
% - Probeer activiteiten in een logische volgorde weer of groepeer ze

% Voor de planning wordt Github gebruikt. Github heeft een ingebouwde mogelijkheid om een roadmap te maken. Er wordt verwacht van de projectleden dat zij 8 uur per week beschikbaar zijn om aan het project te werken, zodat de projectactiviteiten afgerond kunnen worden. 

% Het project bestaat uit een plan van aanpak en drie assesments met bijbehorende activiteiten waarbij op de volgende datum aan wordt gewerkt:
% \begin{itemize}

%   \item \textbf{Plan van aanpak}: 5 februari tot en met 12 februari 
  
%   \item \textbf{Assesment 1}: 13 februari tot en met 25 maart
%   \begin{itemize}
%       \item  Complete the Plan of Approach.
%       \item Design a circuit to read the HAL sensors.
%       \item Design code to read the HAL sensors.
%       \item Design code to control the timing of the MOSFETS 
%       using the HAL sensor data.
%       \item Design code to create an interface using I2C or UART to display data on a touch-display
%       connected to a Raspberry pi model 3B. Or design code to display data on a SSD1306
%       I2C OLED display and to connect a psychical interface using buttons and potentiometer.
%       \item Design a schematic with the most important components to drive the motor and order
%       these parts.
%       \item Document in pdf format listing all of the detailed design requirements
% in relation to the software, including both interface and functionality requirements.
%       \item Well documented and well formatted version of their code in its original file format
%       \item Copy of the same code as above but in .pdf format
%       \item Maximum two page summary of what tests have been completed to
%     date, including any tests on actual prototype hardware in pdf format. 
%     \item Een presentatie geven.
%   \end{itemize}
     
 
     
%   \item \textbf{Assesment 2}: 26 maart tot en met 29 april
%      \begin{itemize}
%         \item Based on testing of the schematic add necessary components, protective circuitry and
%         additional filters to the schematic.
%         \item Test the function of the updated schematic and continue to improve and test the schematic
%         until the schematic functions as is required.
%         \item After completing the schematic the PCB can be designed.
%        \item  Document in pdf format listing all of the detailed design requirements
%         in relation to the hardware (circuits and PCB).
%         \item A zip file containing all PCB related design files (schematic and layout
%         files) as well as datasheets for all active components. This should all
%         be contained in a clearly labeled and organized file structure.
%         \item A well formatted pdf file containing pages showing both the schematic
%         and the layout for the PCB design. The layout should be in a clear
%         format such that it can be easily reviewed both from the top level
%         design perspective and the individual layers.
%         \item Maximum two page summary of the progress of the project, including
%         the results of any tests on actual prototype hardware.
%         \item Een presentatie geven.
%      \end{itemize}
     
%   \item \textbf{Final Assesment}: 30 april tot en met 17 juni
%     \begin{itemize}
%         \item Een presentatie geven
%         \item Een demonstratie geven
%         \item Eindrapport inleveren
%     \end{itemize}
    
% \end{itemize}


\section{Product}

% Beschrijf in dit hoofdstuk welke producten je oplevert waar je met de uitvoering van de activiteiten naar
% toewerkt. Deze producten leiden dan samen naar het eindproduct. Voorbeelden van producten zijn:
% - Een verslag;
% - Tijdsplanning;
% - Een ontwerp
% - Een meetrapport
% - Een tekening;
% - Etc.

During the \textbf{PROJE4} course \textbf{Alset Innovations} will be designing a motor driver to drive a 78W motor continuously. To achieve this a set of deliverables must be completed and delivered. These deliverables are mainly Reports and documents about the progress of the projeect. A detailed list is provided below.

\begin{itemize}
    \item Plan of Approach
        \begin{itemize}
            \item A detailed report about the Plan of Approach of \textbf{Alset Innovations}.
        \end{itemize}
    \item Intermediate Assessment 1
        \begin{itemize}
            \item Document in pdf format listing all of the detailed design requirements
            in relation to the software, including both interface and functionality
            requirements.
            \item Well documented and formatted code in original an pdf file formats.
            \item Summary of completed tests including hardware and software tests.
        \end{itemize}
    \item Intermediate Assessment 2
        \begin{itemize}
            \item Document in pdf format listing all of the detailed design requirements
            in relation to the hardware.
            \item A zip file containing all PCB related design files as well as data-sheets for all active components.
            \item A well formatted pdf file containing pages showing both the schematic
            and the layout for the PCB design. The layout should be in a clear
            format such that it can be easily reviewed both from the top level
            design perspective and the individual layers.
            \item Maximum two page summary of the progress of the project, including
            the results of any tests on actual prototype hardware.
        \end{itemize}
    \item Final Assessment
        \begin{itemize}
            \item A detailed final report according to the provided template. In addition each student is required to write about their competencies.
        \end{itemize}
\end{itemize}
\section{Quality}
% De lezer leest in dit hoofdstuk hoe je de kwaliteit van de (tussen)producten uit de vorige paragraaf
% waarborgt. Dit kun je bijvoorbeeld beschrijven door:
% -
% - Of en in welke mate je gebruikmaakt van extern advies (bijv. iemand buiten je projectgroep die
% je (tussen)product controleert). Geef aan hoe je de feedback verwerkt;
% - Hoe je omgaat met regelgeving, veiligheid, programma van eisen etc;

% Te beschrijven welke software, materialen en andere hulpbronnen je gebruikt en hoet dit
% bijdraagt aan de kwaliteit van jullie projectresultaat.

% Om de kwaliteit voor het project te waarborgen zal er gebruik gemaakt worden van de software STM32CubeIDE, waarmee de stm32-microcontroller geprogrammeerd kan worden. De stm32-chip is gebruikelijk in de industrie, vanwege zijn manier van programmeren, zoals het instellen van de GPIO en klokfrequenties. Bovendien zou het realtime simulatiepakket Caspoc gebruikt kunnen worden. Hiermee kan de werking en aansturing van de motordriver getest worden. Daarnaast wordt KiCad 7.0 gebruikt voor het ontwerpen van PCB's. Kicad is opensource en beschikt over de functies om een PCB te ontwerpen voor de motordriver. Het bestellen van de PCB wordt bij Eurocircuits gedaan, vanwege afspraken die school heeft gemaakt. \\
% Voor het aanschaffen van materialen zijn bedrijven zoals Conrad, Farnell, Texas Instruments en Mouser geschikt, omdat zij beschikken over een aanbod van hardware voor de motordriver.\\
% Voor het vragen van advies kan er terecht bij:
% \begin{itemize}
%   \item Jesse op den Brouw: C-programmeur
%   \item Stephen O'Loughlin: algemeen advies
%   \item Ad van den Bergh: programmeur
%   \item Diëgo Zuidervliet: kennis voor het aansturen van motordrivers
%   \item Paul Witte: PCB-designer
% \end{itemize}
% Voor het controleren van de (tussen)producten worden assesments gegeven met de docenten die hiervoor een aantal punten geven. Bovendien wordt er elke week een vergadering georganiseerd, waarbij de begeleider betrokken is en eventueel feedback kan geven.\\
% Het ontvangen van feedback wordt meegenomen in de toekomstige activiteiten om zo tot een beter ontwerp te komen. Daarnaast hoeft er geen rekening gehouden te worden met de regelgeving, omdat het eindproduct voor eigengebruik is. Voor de veiligheid moeten componenten gekoeld worden en niet kunnen ontploffen. Het eindproduct zal de minimale eisen hebben volgens het programma van eisen, maar de mogelijkheid bestaat dat er extra functies worden toegevoegd.\\
% Alle bovengenoemde producten moeten zorgen dat er een functionele motordriver ontworpen kan worden.



To ensure the quality of the project, the following software and resources will be utilized:

\begin{itemize}
  \item \textbf{STM32CubeIDE}: This software will be used for programming the STM32 microcontroller, which is commonly used in the industry due to its programming capabilities such as configuring GPIO and clock frequencies.
  \item \textbf{Caspoc Real-time Simulation Package}: This tool will be employed to test the operation and control of the motor driver in real-time simulations.
  \item \textbf{KiCad 7.0}: KiCad, an open-source software, will be utilized for designing PCBs for the motor driver, as it provides the necessary features for PCB design.
  \item \textbf{Eurocircuits}: PCB manufacturing will be done through Eurocircuits as per the agreements made by the school.
\end{itemize}

For procurement of materials, companies such as Conrad, Farnell, Texas Instruments, and Mouser will be considered due to their hardware offerings for the motor driver.

Expert advice can be sought from the following individuals:
\begin{itemize}
  \item Jesse op den Brouw: C programmer
  \item Stephen O'Loughlin: General advice
  \item Ad van den Bergh: Programmer
  \item Diëgo Zuidervliet: Expertise in motor driver control
  \item Paul Witte: PCB designer
\end{itemize}

Assessments will be conducted by instructors to evaluate the (intermediate) products. Additionally, weekly meetings will be held involving the supervisor to provide feedback.

Feedback received will be incorporated into future activities to enhance the design. Regulatory considerations are not necessary as the end product is for personal use. However, components must be cooled for safety, and precautions must be taken to prevent explosions. The final product will meet the minimum requirements outlined in the specifications, with the possibility of additional features being added.

All the resources will contribute to the development of a functional motor driver.


\phantomsection
\addcontentsline{toc}{section}{Referenties}
\printbibliography
\end{document}
